
\chapter{INTRODUCTION}
\section{Background}

Vietnam's higher education system has undergone significant changes over the past decade, driven largely by a national push toward digital learning. The government set out concrete targets through Decision No.~1373/QD-TTg~\cite{vietnam_strategy2021}, which mandates the integration of digital technology into teaching at all levels. This goes beyond buying hardware. It reflects a deliberate shift in how the country thinks about educational content creation, distribution, and consumption. Universities affiliated with Vietnam National University -- Ho Chi Minh City (VNU-HCM) are expected to lead that effort.

The VNU-HCM MOOC platform~\cite{vnuhcm_mooc} (\url{https://courses.vnuhcm.edu.vn}) already delivers thousands of course modules and serves as the university's main digital learning infrastructure. On paper, the infrastructure is there. In practice, the gap between having it and actually using it well is wide. Surveys at HCMIU found that 89\% of instructors want to create interactive digital content, but only 23\% feel confident using the available authoring tools. The obstacles they describe consistently come down to three things: interface complexity, steep learning curves, and no real support for instructional design decisions. Studies on e-learning adoption in comparable university contexts show that perceived ease of use is often the deciding factor in whether faculty adopt a tool at all~\cite{hadullo2017}. The bottleneck isn't interest or infrastructure---it's tools that fit how educators actually work.

Interactive video has become one of the most studied formats for asynchronous and hybrid instruction, and the evidence for its effectiveness is solid. The core advantage is straightforward: instead of passively watching a recorded lecture, students must stop and respond to questions embedded at specific timestamps. This forces active retrieval, which the learning science literature consistently supports~\cite{freeman2014}. Merkt et al.\ showed that videos with embedded forced-choice questions produced significantly better knowledge retention and transfer scores compared to equivalent linear video~\cite{merkt2011}. Vural found similar results---strategic timestamp placement improved post-test scores, enhanced long-term retention, and reduced dropout in online courses~\cite{vural2013}.

The design of those embedded questions matters too. Bloom's revised taxonomy~\cite{anderson2001} gives instructors a structured way to think about cognitive levels, from basic recall through to synthesis and evaluation. When questions are placed at the right cognitive level for the concept just covered, and positioned immediately after that concept ends, the result is reduced extraneous load and increased germane load---the kind of mental effort that builds lasting understanding~\cite{sweller1988,chandler1991,paas2003}.

H5P (HTML5 Package) is the open-source framework that makes this kind of interactive video creation technically possible~\cite{h5p2023,lumieducation2023}. Originally developed by Joubel AS in 2013, H5P now covers over 50 content types---from basic multiple choice and true/false to branching scenarios, virtual tours, and interactive timelines. Every H5P package is self-contained: a single \texttt{.h5p} ZIP archive containing its own JavaScript libraries, styling, media files, and a structured JSON configuration. This portability means an H5P package can be imported into any H5P-aware LMS like Moodle~\cite{moodle2023}, Canvas, or Blackboard via standard LTI~1.3~\cite{lti2019}.

H5P Interactive Video is the most widely used content type in higher education. It wraps a video source---either an uploaded file or a YouTube URL---in a player that pauses at pre-defined timestamps to show interactive overlays. Learners must respond before playback resumes. The interaction types most relevant to this project are Multiple Choice, True/False, Fill in the Blanks, Drag Text, Mark the Words, and Matching. Each type maps to a dedicated H5P JavaScript library with a strictly defined JSON schema, which matters for the AI pipeline described later.

Despite H5P's capabilities, the dominant way Vietnamese universities deploy it is through the WordPress H5P plugin---a tool designed for small personal sites, not institutional content production. Creating a single interactive video requires navigating a long sequence of steps: logging into a WordPress admin panel, finding or creating a Post, activating the H5P shortcode block, opening the H5P editor in a separate iframe, uploading the video, manually locating each timestamp through the built-in player, configuring each interaction across several nested settings screens, then saving and returning to the main post. This spans at least six screens and involves web publishing concepts that have nothing to do with teaching. First-time users typically take between 45 and 60 minutes to produce even a basic three-question interactive video.

Research on educational technology adoption points to the Technology Acceptance Model (TAM)~\cite{davis1989} as a useful framework for understanding why educators accept or reject tools. Perceived ease of use and perceived usefulness together account for a large share of adoption variance. When a tool imposes high extraneous cognitive load on every interaction---as the WordPress H5P workflow does through its fragmented interface---perceived ease of use collapses. Adoption fails regardless of how strong the underlying educational evidence might be~\cite{sweller1988,chandler1991}.

Recent advances in large language models have opened up a new possibility for educational content authoring: automatically analyzing a video transcript and generating contextually appropriate, pedagogically structured questions~\cite{brown2020,anthropic2024}. If an instructor can shift from manually creating each question to reviewing and accepting AI-generated suggestions, the cognitive overhead of interactive video production drops substantially. The same pipeline that extracts a transcript, segments it by topic, and returns structured H5P interaction suggestions can also apply Bloom's taxonomy labels to each question---helping instructors spread cognitive difficulty across the video without needing deep instructional design training.

This thesis investigates a specific question: can a purpose-built web authoring interface, combined with an AI pipeline, measurably improve educator productivity and satisfaction when creating H5P interactive videos? The primary AI provider is Groq's fast cloud API using the \texttt{llama-3.3-70b-versatile} model, with a locally hosted Ollama server available as a privacy-preserving fallback.

% ─────────────────────────────────────────────────────────────────────────────
\section{Problem Statement}

The central problem this thesis addresses is the gap between what H5P interactive video can do pedagogically and what instructors can actually accomplish with existing authoring tools.

That gap shows up through four related challenges that reinforce each other:

\begin{enumerate}
  \item \textbf{Workflow complexity and fragmentation.} The WordPress-based H5P workflow requires navigating at least six distinct screens and understanding web publishing concepts that have no connection to the teaching task. A detailed task analysis counted 47 interface interactions---clicks, scrolls, form entries---across those six screens to produce a single interactive video. Page load times in institutional WordPress deployments average 3.2 seconds, which sounds tolerable until you multiply it across 47 sequential interactions: that's nearly two and a half minutes of waiting, while losing the video position each time.

  \item \textbf{High extraneous cognitive load.} The mental effort required to operate the interface leaves less room for thinking about the actual teaching content~\cite{chandler1991,mayer2009}. Instructors in preliminary interviews reported spending more time managing tool settings, timestamp alignments, and shortcode placements than designing questions. That inversion of effort is a sign the tool is working against its users.

  \item \textbf{Low institutional adoption rates.} Because the workflow is this demanding, only faculty with significant IT experience actually use it at VNU-HCM. Instructors in humanities, social science, and other non-technical disciplines---who might benefit most from interactive video---are effectively excluded from a tool with strong empirical evidence behind it~\cite{davis1989,hadullo2017}.

  \item \textbf{No AI assistance.} None of the freely available H5P authoring tools offer integrated AI support for question generation. Instructors without instructional design training have no streamlined way to get context-aware question suggestions tied to what their video is actually about.
\end{enumerate}

The \textbf{primary research question} is:

\begin{quote}
  \itshape How can a purpose-built web platform combining user-centred interface design with an AI-driven transcript analysis pipeline improve the efficiency, pedagogical quality, and institutional adoption of H5P interactive video authoring by university instructors?
\end{quote}

Four supporting questions guide the investigation:

\begin{enumerate}[label=RQ\arabic*.]
  \item What are the specific usability barriers in current H5P authoring workflows that most strongly reduce task completion rates and user satisfaction?
  \item How can user-centred design principles be applied to reduce extraneous cognitive load in timestamp-based interactive question authoring?
  \item What full-stack architecture best supports scalable, real-time AI-assisted question generation and H5P package assembly in a unified web environment?
  \item To what extent does the newly developed platform quantitatively and qualitatively outperform the WordPress baseline on usability, efficiency, and productivity metrics?
\end{enumerate}

% ─────────────────────────────────────────────────────────────────────────────
\section{Scope and Objectives}

The thesis pursues five concrete objectives to answer these research questions.

\begin{enumerate}
  \item \textbf{Platform Development and Architecture.} Design and build a full-stack web application that lets instructors upload or link a video, add H5P interactions at specific timestamps through a visual timeline interface, preview content in real time, and export a standards-compliant \texttt{.h5p} package ready for import into the VNU-HCM MOOC infrastructure and other LTI-compatible LMS environments.

  \item \textbf{Advanced AI Pipeline Integration.} Build and integrate a multi-provider AI pipeline that extracts timestamped transcripts from uploaded videos (via YouTube captions, uploaded WebVTT/SRT files, or local Whisper transcription), sends the transcript to a large language model, and returns a validated JSON payload containing H5P interaction suggestions. Each suggestion should include precise timestamps, the appropriate question type, distractor options, and a Bloom's taxonomy classification.

  \item \textbf{User-Centred Design.} Apply user-centred design methodology throughout the interface design and prototyping process. All design decisions must be grounded in identified instructor needs and validated against Nielsen's usability heuristics.

  \item \textbf{Security Implementation.} Implement a security baseline aligned with the OWASP Top Ten, including JWT-based authentication, role-based access control, server-side Zod input validation, HTTPS enforcement, and rate limiting.

  \item \textbf{Empirical Usability Evaluation.} Conduct a controlled, comparative usability study measuring task completion rate, time-on-task, error frequency, and SUS scores against the WordPress H5P baseline. Supplement quantitative data with think-aloud protocols and post-task interviews.
\end{enumerate}

% ─────────────────────────────────────────────────────────────────────────────
The scope is bounded by the following constraints.

\textbf{Content type scope.} The platform specifically targets six H5P interaction types: Multiple Choice, True/False, Fill in the Blanks, Drag Text, Mark the Words, and Matching. All six are supported for manual authoring. The AI pipeline generates all six types based on a pattern matrix that maps content characteristics to question types. These types have well-defined JSON schemas that can be reliably validated with Zod, ensuring AI output never breaks the H5P player.

\textbf{Video source scope.} The platform accepts direct video file uploads up to 500~MB (MP4, WebM, MOV) and YouTube video URLs. YouTube-sourced transcripts are extracted from the platform's automatic or manual captions API. For uploaded files, instructors can provide a WebVTT caption file or let the system invoke \texttt{faster-whisper} locally for automatic audio transcription.

\textbf{User scope.} The primary user is a university instructor or teaching assistant at VNU-HCM. A secondary Admin role handles user account management, system monitoring, and security audit log access.

\textbf{Integration scope.} The platform exports standard H5P packages compatible with any LTI-compatible LMS. An LTI~1.3 launch route is stubbed in the backend architecture, but LMS gradebook passback and xAPI learning analytics are designated as future work.

Several limitations are acknowledged upfront:

\begin{itemize}
  \item \textbf{H5P content types.} Complex types like Branching Scenarios, Course Presentations, and Virtual Tours are out of scope.
  \item \textbf{Evaluation sample.} Usability testing used ten senior IT students as educator proxies. This is standard practice for early-stage formative evaluation~\cite{nielsen1994}, but it limits generalizability to less technically experienced demographics.
  \item \textbf{Mobile authoring.} The platform is designed for desktop and large tablet browsers. Mobile authoring is technically possible but was not a primary design goal and was not formally evaluated.
  \item \textbf{Concurrent load.} The system was stress-tested at up to 50 concurrent users. Large-scale institutional deployment would require a more robust architecture, likely Kubernetes-based.
  \item \textbf{AI transcription quality.} The quality of AI-generated suggestions depends directly on transcript accuracy. Poor audio quality or heavy accents will degrade Whisper output and, consequently, the AI suggestions.
\end{itemize}

% ─────────────────────────────────────────────────────────────────────────────
\section{Research Contributions}

This thesis makes five contributions to educational technology research and software engineering practice.

\begin{enumerate}
  \item \textbf{A reproducible platform architecture.} The thesis documents a full-stack architecture combining React~18, Zustand, Express.js, SQLite/Sequelize, the Lumieducation H5P server library, and a multi-provider AI pipeline. The documentation is detailed enough for other universities to reproduce, modify, and deploy.
  \item \textbf{A novel AI-assisted authoring workflow.} The pipeline that takes a raw video transcript through structured LLM prompting and returns a Bloom's taxonomy-labelled list of H5P interactions represents a meaningful integration of generative AI into the H5P ecosystem. The two-step prompting approach---topic segmentation first, then question generation per topic using a content-type pattern matrix---is specifically designed to reduce hallucination and produce structured output.
  \item \textbf{Prompt engineering for structured H5P JSON output.} The thesis details the system prompts, few-shot examples, and Zod validation strategies developed to reliably coerce an LLM into returning complex JSON payloads that conform to H5P content type specifications. This provides a reusable design pattern for other structured AI output scenarios.
  \item \textbf{Empirical usability evidence.} The comparative evaluation against the WordPress H5P baseline provides quantitative evidence that a purpose-built, user-centred platform can substantially reduce authoring time, improve task completion rates, and raise usability scores on the System Usability Scale.
  \item \textbf{Open-source contribution.} The full codebase---including AI service routing, prompt templates, database schemas, and H5P integration middleware---is made publicly available to support further research and free institutional deployment.
\end{enumerate}

The practical value runs in three directions. By lowering the technical barrier, the platform enables a wider range of VNU-HCM faculty to independently produce interactive content. The AI-assisted workflow demonstrates that the gap between generative AI research and reliable educational tool development can be bridged with careful software engineering. And the platform aligns with Vietnam's national digital transformation goals~\cite{vietnam_strategy2021} by providing a locally deployable, open-source solution that reduces reliance on expensive proprietary software.

% ─────────────────────────────────────────────────────────────────────────────
\section{Structure of thesis}

The rest of this report is organized as follows.

\textbf{Chapter~2 -- Literature Review} covers the theoretical and technical foundations of the work. It examines the state of e-learning in Vietnamese higher education, the H5P framework and its architectural limitations, empirical research on interactive video pedagogy, cognitive load theory, user-centred design principles, the evolution of large language models, prompt engineering techniques, and LMS integration standards. The chapter ends with a critical survey of related systems and an identification of the specific research gaps this project addresses.

\textbf{Chapter~3 -- Methodology} describes the research design, requirements analysis process, user-centred design methodology, agile development approach, and empirical evaluation protocol. It also provides a detailed technical description of the AI pipeline architecture, including the prompt engineering strategy and pattern matrix used for question type selection.

\textbf{Chapter~4 -- Implementation and Results} covers the complete system implementation and measured technical outcomes. It details the backend and frontend architecture, the H5P package assembly mechanism, the AI suggestion injection pipeline, security controls, and performance measurements including AI latency and API load results. Scalability analysis and Vietnamese language support are also covered here.

\textbf{Chapter~5 -- Discussion and Evaluation} presents the results of the comparative usability study, including task completion rate, time-on-task, and SUS scores, and interprets the qualitative themes from participant feedback.

\textbf{Chapter~6 -- Conclusion and Future Work} synthesizes the findings, assesses the project's contributions against the research questions, discusses limitations, and proposes a roadmap for future development.
